\scanpage{199}\mspage{- 58 -}\hfill IX

qui est une \uline{dualité parfaite}, i.e. identifiant chacun des deux facteurs du premier membre au dual de l'autre (au sens de 5.2.2). (NB les exposants ab désignent les ``parties abéliennes'', quotient de la partie fixe par la partie torique). Il est clair d'autre part que (5.3.2) est induit par l'accouplement

\[
T_{\handwritten{\ell}}(A^{\unsure{o|0}})^{\mathrm{ab}}\times T_{\handwritten{\ell}}(A^{\prime o})^{\mathrm{ab}}\handwritten{\longrightarrow} T_{\handwritten{\ell}}(\mathbb{G}_{mS})\tag*{(5.3.3)}
\]

déduit par passages aux quotients de l'accouplement (5.2.4.2). Celui-ci étant une dualité parfaite sur les fibres génériques, i.e. l'homomorphisme correspondant

\[
T_{\handwritten{\ell}}(A^o)^{\mathrm{ab}}\handwritten{\longrightarrow} D(T_{\handwritten{\ell}}(A^{\prime o})^{\mathrm{ab}})\tag*{(5.3.4)}
\]

étant un isomorphisme sur les fibres génériques, nous avons évidemment envie d'en déduire qu'il en est de même sur $S$ tout entier, ou ce qui revient au même, sur les fibres spéciales. Cette déduction est valable, trivialement, si $\handwritten{\ell}\handwritten{\ne}p$ ; elle est valable également, du moins dans le cas d'inégales caractéristiques, si $\handwritten{\ell}=p$, grâce à un théorème profond de TATE [33, th.4]. Nous allons cependant pouvoir l'établir directement, sans l'aide du théorème de Tate, plus bas (7.4.6).

Pour préciser la signification géométrique du résultat que nous désirons établir, notons que l'accouplement déduit de (5.3.3) par passage aux fibres spéciales n'est autre que l'accouplement

\[
T_{\handwritten{\ell}}(B_o)\times T_{\handwritten{\ell}}(B'_o)\handwritten{\longrightarrow} T_{\handwritten{\ell}}(\mathbb{G}_{mk})\tag*{(5.3.5)}
\]

déduit de la biextension $W_o$ de $(A_o^o,A^{\prime o}_o)$ par passage au quotient, $B_o$ et $B'_o$

% archive sheet number: 198
% archive numbering: 376 (pencil and printed label, foot).
